\input{../tex-inputs/latex-leseansicht-vorspann}

\section[Gustav Schwarzkopf an Arthur Schnitzler, 4. 2. 1901]{L04306 Gustav Schwarzkopf an Arthur Schnitzler, 4. 2. 1901}
\nopagebreak\mylabel{L04306v}
\rehead{ }\normalsize\beginnumbering\briefempfaengerindex{Schnitzler, Arthur@\textsc{Schnitzler, Arthur}!zzzSchwarzkopf, Gustav@\emph{von Gustav Schwarzkopf}!1901-02-041@{4. 2. 1901}|(be}
\toendnotes[C]{\smallbreak\pagebreak[2]}
\correspDesc{Versand  durch Gustav Schwarzkopf am 4. 2. 1901 in Wien
\newline{}Übermittlung  am 4. 2. 1901 in Wien
\newline{}Zustellung  am 5. 2. 1901 in Wien
\newline{}Erhalt  durch Arthur Schnitzler am 5. 2. 1901 in Wien}\toendnotes[C]{\smallbreak}
\Standort{CUL, Schnitzler, B 96a.}
\physDesc{Postkarte, 361 Zeichen
\newline{}Handschrift: schwarze Tinte, deutsche Kurrent
\newline{}Versand: 1) Stempel: »\nobreak{}\oindex{I., Innere Stadt@\textbf{I., Innere Stadt}, \emph{Verwaltungsgebiet}|pwk}1/1 Wien 1, 4. 2. 01, 7–8N\nobreak{}«.   2) Stempel: »\nobreak{}\oindex{Wien@\textbf{Wien}, \emph{Verwaltungsgebiet}|pwk}Wien, 5. 2. \textcolor{gray}{01}, 8\nobreak{}«. }\toendnotes[C]{\smallbreak}\pstart{}{\pb}Herrn \textsc{D\textsuperscript{r} Arthur Schnitzler}\pend{}\pstart{}\textsc{Wien\oindex{Wien@\textbf{Wien}, \emph{Verwaltungsgebiet}|pw}}\pend{}\pstart{}IX. Frankgaſſe 1\oindex{Wien@\textbf{Wien}!IX., Alsergrund@\textbf{IX., Alsergrund}!Frankgasse 1@\textbf{Frankgasse 1}, \emph{Wohngebäude}|pw}.\pend{}\pstart{}II. Stock\pend{}{\bigskip}\vspace{1em}
\pstart
           \raggedleft{}{\pb}4. II 1901\pend
           \vspace{0.5em}
\pstart
           Lieber Arthur! Beſten Dank, aber es wird{ }ſich leider nicht machen
               laſſen. Ich bin nicht ganz wol, – wahrſcheinlich eine leichte Influenza – und dürfte
                  \label{K_L04306-1v}\edtext{morgen}{\lemma{\textnormal{\emph{morgen}}}\Cendnote{\textnormal{Am 5. 2. 1901 fand bei Schnitzler
                  in der Frankgasse 1\oindex{Wien@\textbf{Wien}!IX., Alsergrund@\textbf{IX., Alsergrund}!Frankgasse 1@\textbf{Frankgasse 1}, \emph{Wohngebäude}|pwk} eine private Lesung\eventindex{Frankgasse 1@\textbf{Frankgasse 1}!Private Lesung von Die Gedenktafel der Prinzessin Anna, Das Haus Delorme, Marionetten, 5.2.1901@Private Lesung von Die Gedenktafel der Prinzessin Anna, Das Haus Delorme, Marionetten, 5.2.1901|pwkv} statt, wobei nur durch
                  das vorliegende Korrespondenzstück die betreffende Stelle im Eintrag zum 5. 2. 1901 im \emph{Tagebuch}\pwindex{Schnitzler, Arthur 15.\,5.\,1862 Wien – 21.\,10.\,1931 ebd.@\textsc{Schnitzler, Arthur} (15.\,5.\,1862 Wien – 21.\,10.\,1931 ebd.), \emph{Schriftsteller, Mediziner}!Tagebuch@\strich\emph{Tagebuch}|pwk} zu begreifen ist. Felix Salten\pwindex{Salten, Felix 6.\,9.\,1869 Budapest – 8.\,10.\,1945 Zürich@\textsc{Salten, Felix} (6.\,9.\,1869 Budapest – 8.\,10.\,1945 Zürich), \emph{Schriftsteller, Journalist, Chefredakteur}|pwk} las \emph{Die
                     Gedenktafel der Prinzessin Anna}\pwindex{Salten, Felix 6.\,9.\,1869 Budapest – 8.\,10.\,1945 Zürich@\textsc{Salten, Felix} (6.\,9.\,1869 Budapest – 8.\,10.\,1945 Zürich), \emph{Schriftsteller, Journalist, Chefredakteur}!Gedenktafel der Prinzessin Anna@\strich\emph{Die Gedenktafel der Prinzessin Anna}|pwk}, Schnitzler{ }\emph{Das Haus Delorme}\pwindex{Schnitzler, Arthur 15.\,5.\,1862 Wien – 21.\,10.\,1931 ebd.@\textsc{Schnitzler, Arthur} (15.\,5.\,1862 Wien – 21.\,10.\,1931 ebd.), \emph{Schriftsteller, Mediziner}!Haus Delorme. Eine Familienszene@\strich\emph{Das Haus Delorme. Eine Familienszene}|pwk} und \emph{Marionetten}\pwindex{Schnitzler, Arthur 15.\,5.\,1862 Wien – 21.\,10.\,1931 ebd.@\textsc{Schnitzler, Arthur} (15.\,5.\,1862 Wien – 21.\,10.\,1931 ebd.), \emph{Schriftsteller, Mediziner}!Zum großen Wurstel. Burleske in einem Akt@\strich\emph{Zum großen Wurstel. Burleske in einem Akt}|pwk} (späterer Titel: \emph{Zum großen
                     Wurstel}\pwindex{Schnitzler, Arthur 15.\,5.\,1862 Wien – 21.\,10.\,1931 ebd.@\textsc{Schnitzler, Arthur} (15.\,5.\,1862 Wien – 21.\,10.\,1931 ebd.), \emph{Schriftsteller, Mediziner}!Zum großen Wurstel. Burleske in einem Akt@\strich\emph{Zum großen Wurstel. Burleske in einem Akt}|pwk}). Weitere Anwesende sind nicht dokumentiert. }}}\label{K_L04306-1} noch nicht
               geſellſschaftsfähig{ }ſein. Ich bitte Sie Salten\pwindex{Salten, Felix 6.\,9.\,1869 Budapest – 8.\,10.\,1945 Zürich@\textsc{Salten, Felix} (6.\,9.\,1869 Budapest – 8.\,10.\,1945 Zürich), \emph{Schriftsteller, Journalist, Chefredakteur}|pw}
               mein Bedauern auszuſprechen und mich bei ihm zu entſchuldigen.\pend
           
\pstart
           Herzlichſt{\\[\baselineskip]}Ihr{\\[\baselineskip]}\spacefill\mbox{Gustav Sch.}\pend
           \leftskip=0em{}\selectlanguage{ngerman}\endnumbering\briefempfaengerindex{Schnitzler, Arthur@\textsc{Schnitzler, Arthur}!zzzSchwarzkopf, Gustav@\emph{von Gustav Schwarzkopf}!1901-02-041@{4. 2. 1901}|)be}\mylabel{L04306h}
\begin{anhang}
\end{anhang}\newcommand{\dateiname}{L04306}\newcommand{\titel}{Gustav Schwarzkopf an Arthur Schnitzler, 4. 2. 1901}\newcommand{\editorInnen}{Herausgegeben von Jahnke, SelmaMüller, Martin Anton}\input{../tex-inputs/latex-leseansicht-abspann}
